% ===== §0 The Delineation of the Phenomenon =====
\section*{The Phenomenon}
\addcontentsline{toc}{section}{The Phenomenon}
\label{sec:phenomenon}

\noindent
Before any theory is brought to bear, the object to be explained must be set out in its own terms. This preliminary section does only that. It describes a phenomenon and marks it off from its neighbours, using as little theoretical vocabulary as possible; the terms by which the later chapters will explain it---``generativity,'' ``the symbolic,'' and the rest---are deliberately withheld here, since to describe the phenomenon in the language of one's eventual explanation is already to prejudge it. What follows is a description of how the thing is given, not an account of what it is.

\subsection*{The datum}
\label{sub:phen-datum}

Consider two people who, on an afternoon with nothing to do, sit together and watch the rain beyond a window. They are not talking, or not much; there is nothing they are trying to settle, nothing they are waiting for, no task to complete and no decision to reach. Nothing, in any ordinary sense, is happening. And yet they are happy---not euphoric, not moved to any peak, but plainly and unmistakably content, in a way neither would hesitate to call happiness if asked.

The phenomenon has a feature that makes it difficult to place: it is full without being eventful. The contentment is not the residue of something that has just occurred, nor the anticipation of something about to occur. It does not point beyond the afternoon to an achievement secured or a lack supplied. It is simply, and without remainder, the quiet fullness of the present hour. This is what the standard theories of happiness find hard to accommodate, and what the rest of this paper is written to address. For the moment it is enough to record it as a datum: there is a happiness that consists in nothing happening, and that is nonetheless full.

\subsection*{The features of the phenomenon}
\label{sub:phen-features}

Stated more carefully, the phenomenon shows several features, each of which can be recognised without any theory.

\emph{Uneventfulness.} Nothing is achieved, undergone, or resolved. No new state is reached and no event occurs. This is the feature that most resists the received accounts, which tend to locate happiness in an attainment or in the intensity of an experience.

\emph{Fullness.} Despite the absence of event, the condition is not one of boredom, emptiness, or torpor. It is full. Here ``full'' is used only as a plain description of how the condition is felt, not as a theoretical claim; the question of what this fullness consists in is left entirely open.

\emph{Sharedness.} The happiness is essentially a happiness of being \emph{together}. It is not felt as belonging to either party alone but as belonging to the situation the two of them are in. A single person reproducing the same outward scene---the same window, the same rain---would not necessarily be in the same condition. Whatever the fullness is, it is located in the being-with, not in either of the two.

\emph{A faint movement, not inertness.} The condition is still, but not with the stillness of a stone. There is in it something difficult to name---a slight, unhurried liveliness: the shared breathing, the common direction of attention toward the same rain, the slow passage of the hour. It is not a stop. It is more like a quiet, unhurried proceeding that is not in any haste to arrive anywhere.

\emph{Absence of want.} There is no sense of lack, no felt pressure of ``something still missing.'' And just for this reason there is nothing that needs to be said and nothing that needs to be settled---a plain, unspoken ease.

\subsection*{Demarcation from neighbouring conditions}
\label{sub:phen-demarcation}

The phenomenon is easily confused with several conditions it resembles, and the description is not complete until these are set apart. The distinctions can be felt at the descriptive level, before any theory; their theoretical characterisation is deferred.

\emph{Against boredom and depressive withdrawal.} These too are ``uneventful,'' but they are conditions of staleness, emptiness, or entrapment, whereas the present phenomenon is one of ease and fullness. The two share the absence of event and differ in everything else. This already marks something the later chapters must respect: there are two stillnesses, an empty one and a full one, and they are not told apart by whether anything is happening.

\emph{Against satiety or the relaxation of a satisfied need.} Satiety is the aftermath of an event---the having-eaten, the having-obtained---and points back to a want just supplied. The present phenomenon points to no such want and is not the ``after'' of anything. It is not a discharge of tension but a condition with no tension to discharge.

\emph{Against the peak experience.} Rapture or ecstasy is the summit of an intensity and is itself an event. The present phenomenon is the opposite in form: low in intensity, without event, and indefinitely sustainable. It is not a height but a level.

\subsection*{Waiting: a neighbouring mode of the same phenomenon}
\label{sub:phen-waiting}

There is a second scene that belongs to the same family. Two people slow down for one another, wait together for a while, watch the light fade as evening comes on. Here too there are uneventfulness, sharedness, faint movement, and the absence of want; but there is in addition an \emph{orientation}---toward something not yet arrived: the evening coming down, the other catching up.

A plain observation distinguishes this from its counterfeit. This kind of waiting is not restless; it is, if anything, sweet and full. It is wholly unlike the waiting of one who has missed a train---which is fretful, out of one's control, and oriented toward a lack. The two kinds of waiting are felt to be quite different, and the difference can be marked here even though its mechanism and structure are left for later. What this mode shows is that the family of the phenomenon includes not only an ``already at rest'' but also a ``slowly under way''; its modes are plural.

\subsection*{The limit of description, and what follows}
\label{sub:phen-limit}

This section has done only what description can do: it has identified a phenomenon, set out its features, and marked it off from its neighbours. It has not answered---because description cannot answer---what this fullness is, why it is full rather than empty, or where exactly the line falls between the two stillnesses and the two kinds of waiting. These are questions of explanation, not of description, and they require theory.

The order of what follows is accordingly this. The next chapter listens to how the existing theories explain this already-delineated phenomenon, and finds that they converge on a common omission. The chapter after that advances the paper's own account and grounds it. Only then are the modes redrawn, and only then are the several illustrative registers---dynamical, neural---allowed to cast their partial light. The phenomenon, for now, is simply on the table.
